\documentclass[conference,compsoc]{IEEEtran}

\usepackage{amsmath,amsfonts,amssymb}
\usepackage{algorithmic}
\usepackage{algorithm}
\usepackage{array}
\usepackage[caption=false,font=footnotesize,labelfont=sf,textfont=sf]{subfig}
\usepackage{textcomp}
\usepackage{stfloats}
\usepackage{url}
\usepackage{verbatim}
\usepackage{graphicx}
\usepackage[nocompress]{cite}
\usepackage{listings}
\usepackage{xcolor}
\usepackage[catalan]{babel}
\usepackage[T1]{fontenc}
\usepackage[utf8]{inputenc}
\usepackage{booktabs}
\usepackage{float}
\usepackage{cuted}
\usepackage{microtype}
\graphicspath{{./}}

\renewcommand{\abstractname}{Resum}
\renewcommand{\IEEEkeywordsname}{Termes d'índex}

\makeatletter
\def\abstract{\normalfont\@IEEEtweakunitybaselinestretch{1.15}%
    \if@twocolumn
      \@IEEEabskeysecsize\noindent\textbf{\textit{\abstractname}}---\normalfont\@IEEEabskeysecsize\relax
    \else
      \bgroup\par\addvspace{0.5\baselineskip}\centering\vspace{-1.78ex}\@IEEEabskeysecsize\textbf{\abstractname}\par\addvspace{0.5\baselineskip}\egroup\quotation\normalfont\@IEEEabskeysecsize%
    \fi\@IEEEgobbleleadPARNLSP}
\def\IEEEkeywords{\normalfont\@IEEEtweakunitybaselinestretch{1.15}%
    \if@twocolumn
      \@IEEEabskeysecsize\vskip 0.5\baselineskip plus 0.25\baselineskip minus 0.25\baselineskip\noindent
      \textbf{\textit{\IEEEkeywordsname}}---\normalfont\@IEEEabskeysecsize\relax
    \else
      \bgroup\par\addvspace{0.5\baselineskip}\centering\@IEEEabskeysecsize\textbf{\IEEEkeywordsname}\par\addvspace{0.5\baselineskip}\egroup\quotation\normalfont\@IEEEabskeysecsize%
    \fi\@IEEEgobbleleadPARNLSP}
\makeatother

\lstset{
    language=Java,
    basicstyle=\ttfamily\footnotesize,
    keywordstyle=\color{blue}\bfseries,
    commentstyle=\color{gray}\itshape,
    stringstyle=\color{red!70!black},
    numbers=left,
    numberstyle=\tiny\color{gray},
    numbersep=5pt,
    frame=single,
    breaklines=true,
    showstringspaces=false,
    tabsize=4,
    captionpos=b
}

\hyphenation{mi-ni-max heu-ris-ti-ca con-nec-ta gra-ve-tat con-tro-la-dor}

\begin{document}

\title{Agent per a una Variant de Connecta 4\\Pràctica 6}

\author{\IEEEauthorblockN{Serafí Nebot Ginard}
\IEEEauthorblockA{Universitat de les Illes Balears (UIB)\\
Grau en Enginyeria Informàtica\\
Algorismes Avançats, Curs 2025--2026}}

\maketitle

\begin{strip}
\begin{abstract}
Aquesta memòria descriu el disseny i la implementació d'una aplicació Java per
jugar a una variant de Connecta 4 sobre un tauler de $7\times7$. A més de la
inserció vertical clàssica, el joc permet eliminar una fitxa pròpia i rotar el
tauler 90 graus a l'esquerra o a la dreta, aplicant gravetat després de cada
canvi d'orientació. L'aplicació s'ha implementat amb arquitectura
Model-Vista-Controlador, interfície gràfica Swing, animacions de caiguda i
rotació, modes humà contra humà, humà contra agent i agent contra agent, i
reproducció de partides amb historial. S'han desenvolupat tres agents: un agent
aleatori, un agent greedy d'una sola capa i un agent minimax amb profunditat
limitada, ordenació de moviments i poda alfa-beta. La memòria formalitza les
regles implementades, descriu l'espai d'estats i la funció heurística emprada pel
minimax, i discuteix els resultats funcionals obtinguts sense introduir dades
experimentals no mesurades.
\end{abstract}

\begin{IEEEkeywords}
Connecta 4, minimax, poda alfa-beta, branch and bound, heurística, MVC, Java
Swing, jocs adversarials.
\end{IEEEkeywords}
\end{strip}

\IEEEpeerreviewmaketitle

%% ================================================================
\section{Introducció}

El Connecta 4 clàssic és un joc adversarial de dos jugadors en què les fitxes es
col·loquen en columnes i cauen fins a la primera posició lliure. Guanya el primer
jugador que forma quatre fitxes consecutives en línia horitzontal, vertical o
diagonal. La pràctica modifica aquest joc en tres aspectes importants: el tauler
és de $7\times7$, cada jugador pot eliminar una fitxa pròpia i el tauler complet
es pot rotar 90 graus a l'esquerra o a la dreta. Després d'una eliminació o una
rotació, les fitxes tornen a caure segons la nova orientació del tauler.

Aquestes modificacions fan que el problema sigui més interessant que el Connecta
4 estàndard. Una jugada pot afavorir el jugador que mou, però també pot crear una
línia guanyadora per al contrari; fins i tot pot generar quatre en línia per als
dos jugadors alhora, cas que l'enunciat defineix com empat. A més, l'existència
de rotacions i eliminacions fa que el joc no sigui estrictament progressiu en el
mateix sentit que el joc clàssic: el nombre de fitxes pot disminuir i la forma
del tauler pot canviar completament amb una sola acció.

L'objectiu obligatori de la pràctica és construir una aplicació amb interfície
gràfica que permeti jugar contra una màquina i que aquesta determini la jugada
mitjançant Branch and Bound. En aquesta implementació s'ha modelat aquest requisit
com un agent \texttt{MinimaxAgent} amb poda alfa-beta. En jocs de dos jugadors,
la cerca minimax proporciona el marc natural per alternar decisions entre
maximitzador i minimitzador~\cite{ref_gfg_minimax}, mentre que la poda alfa-beta
actua com el mecanisme de Branch and Bound: descarta branques que ja no poden
millorar la decisió final~\cite{ref_gfg_alphabeta}.

Com a millores opcionals s'han incorporat modes agent contra agent, un agent
aleatori i un agent greedy. Això permet observar estratègies de complexitat molt
diferent: una sense coneixement del joc, una heurística local i una cerca
adversarial amb profunditat limitada. La interfície també permet reproduir
partides automàtiques amb retard, aturar-les, avançar jugades i retrocedir dins
de l'historial.

%% ================================================================
\section{Descripció del Problema}

\subsection{Tauler i Torns}

El tauler té $N=7$ files i $N=7$ columnes. Cada cel·la pot estar buida, ocupada
pel jugador vermell o ocupada pel jugador groc. Els jugadors alternen torns, i en
cada torn el jugador actual tria exactament una acció legal. La implementació fa
servir el jugador vermell com a primer jugador.

Les accions possibles són:

\begin{itemize}
    \item \textbf{Inserir}: triar una columna que no estigui plena i col·locar-hi
          una fitxa pròpia. La fitxa cau fins al fons o fins quedar damunt una
          altra fitxa.
    \item \textbf{Eliminar}: triar una fitxa pròpia del tauler i retirar-la.
          Les fitxes que hi havia per damunt cauen per ocupar l'espai.
    \item \textbf{Rotar a l'esquerra}: girar tot el tauler 90 graus en sentit
          antihorari i aplicar gravetat.
    \item \textbf{Rotar a la dreta}: girar tot el tauler 90 graus en sentit
          horari i aplicar gravetat.
\end{itemize}

Un moviment és terminal si, després d'aplicar-lo, hi ha quatre fitxes consecutives
d'algun jugador. Si només un jugador té quatre en línia, aquest jugador guanya,
encara que no sigui necessàriament qui ha fet la jugada. Si tots dos jugadors
tenen quatre en línia després del mateix moviment, la partida acaba en empat. Si
el tauler queda ple i no hi ha cap guanyador, també acaba en empat. Seguint
l'enunciat, quan el tauler és ple no es continua jugant a base d'eliminar fitxes.

\subsection{Espai de Moviments}

En una posició no terminal, el nombre de moviments legals pot ser molt superior
al Connecta 4 clàssic. En el joc original només hi pot haver fins a 7 moviments,
un per columna. En aquesta variant hi pot haver fins a 7 insercions, fins a $p$
eliminacions, on $p$ és el nombre de fitxes pròpies del jugador actual, i dues
rotacions. Per tant, el factor de ramificació màxim és

\begin{equation}
    b \leq 7 + p + 2.
\end{equation}

Com que $p$ pot arribar a ser aproximadament la meitat del tauler, la cerca
exhaustiva completa és massa costosa per fer-la a cada torn dins d'una interfície
interactiva. Per aquest motiu el minimax implementat és de profunditat limitada i
utilitza una funció heurística per valorar les fulles no terminals.

%% ================================================================
\section{Fonaments Teòrics}

\subsection{Model d'Estat}

L'estat del joc es representa com

\begin{equation}
    S = (B, j, r),
\end{equation}

on $B$ és el tauler, $j$ és el jugador a qui toca moure i $r$ és el resultat
actual: partida en curs, victòria del vermell, victòria del groc o empat. Aquesta
separació és la que apareix a la classe \texttt{GameState}: el tauler conté la
geometria i la gravetat, mentre que l'estat decideix quins moviments són legals i
quan s'ha acabat la partida.

La transició d'estat es defineix com

\begin{equation}
    S' = \operatorname{apply}(S, m),
\end{equation}

on $m$ és un moviment legal. Primer es construeix un nou tauler aplicant la física
del moviment, després s'avaluen les condicions de victòria i finalment es canvia
el torn. La implementació utilitza objectes immutables: aplicar un moviment no
modifica l'estat anterior, sinó que retorna un estat nou. Això és especialment
important per al minimax, ja que la cerca explora moltes línies hipotètiques i no
ha de corrompre la partida real.

\subsection{Funció Heurística}

La funció heurística de \texttt{MinimaxAgent} combina quatre components:

\begin{equation}
\begin{aligned}
h(S)=&h_{terminal}(S)+h_{finestres}(S)\\
     &+h_{centre}(S)+h_{material}(S).
\end{aligned}
\end{equation}

El component terminal domina tots els altres. Una victòria del maximitzador rep
un valor molt alt i una derrota rep un valor molt baix. També s'afegeix un petit
ajust segons la profunditat restant: una victòria trobada abans és preferible a
una victòria més tardana, i una derrota tardana és menys dolenta que una derrota
immediata.

El component principal és $h_{finestres}$. Una finestra és qualsevol segment de
quatre cel·les consecutives en una de les quatre direccions guanyadores:
horitzontal, vertical, diagonal principal o diagonal secundària. Cada finestra es
classifica segons quantes fitxes pròpies i contràries conté. Si conté fitxes dels
dos jugadors, queda bloquejada i val 0, perquè cap jugador pot completar exactament
aquella línia. Si només conté fitxes d'un jugador i buits, es valora segons la
proximitat a una línia de quatre.

\begin{table}[H]
\centering
\caption{Pesos emprats per valorar finestres de quatre cel·les.}
\label{tab:window-weights}
\begin{tabular}{lrr}
\toprule
Situació & Punts propis & Punts oponent \\
\midrule
Quatre fitxes & $+100000$ & $-100000$ \\
Tres fitxes & $+900$ & $-1100$ \\
Dues fitxes & $+90$ & $-100$ \\
Una fitxa & $+8$ & $-8$ \\
Bloquejada & $0$ & $0$ \\
\bottomrule
\end{tabular}
\end{table}

La penalització d'una triple amenaça de l'oponent és una mica més forta que la
recompensa d'una triple pròpia. Aquesta decisió fa que l'agent prefereixi bloquejar
perills immediats abans de crear amenaces simètriques que podrien arribar tard.
El control del centre afegeix pes a les fitxes situades a la columna central,
perquè participen en més línies possibles que les columnes extremes. Finalment,
el material és un terme petit que suma les fitxes pròpies i resta les contràries;
no decideix la partida per si sol, però serveix com a desempat quan dues posicions
tenen amenaces similars.

\subsection{Minimax}

El minimax és l'algorisme clàssic per a jocs deterministes de suma zero amb dos
jugadors i informació perfecta~\cite{ref_gfg_minimax}. Des del punt de vista del jugador
que mou, la seva decisió intenta maximitzar el valor de la posició, mentre que el
contrari intenta minimitzar-lo. Per a una profunditat restant $d$, la recurrència
és

\begin{equation}
\begin{aligned}
V(S,d)&=h(S) && \text{si } S \text{ es fulla}\\
V(S,d)&=\max_m V(S_m,d-1) && \text{si } j=j_{max}\\
V(S,d)&=\min_m V(S_m,d-1) && \text{altrament}
\end{aligned}
\end{equation}

on $m$ recorre $M(S)$, $S_m=\operatorname{apply}(S,m)$ i una fulla és un estat
terminal o un estat amb $d=0$. El valor retornat per la cerca no és una
probabilitat, sinó una utilitat relativa: valors positius afavoreixen el jugador
maximitzador i valors negatius afavoreixen l'oponent.

\subsection{Poda Alfa-Beta com a Branch and Bound}

La poda alfa-beta manté dos límits durant la cerca: $\alpha$, el millor valor que
el maximitzador ja pot garantir, i $\beta$, el millor valor que el minimitzador ja
pot garantir~\cite{ref_gfg_alphabeta}. Quan es compleix $\alpha \geq \beta$, la branca actual no pot canviar
la decisió presa en un nivell superior i es pot descartar. Aquesta és la part de
Branch and Bound de la solució: s'explora l'arbre de joc, però es tallen branques
que queden fora dels límits útils.

L'ordenació de moviments no canvia el resultat del minimax, però sí que afecta
molt el nombre de podes. Si els moviments bons es proven primer, els límits
$\alpha$ i $\beta$ es tanquen abans i la cerca evita més successors. Per això
\texttt{MinimaxAgent} ordena primer moviments amb victòria immediata i després
prefereix insercions centrals.


%% ================================================================
\section{Disseny i Implementació}

\subsection{Arquitectura MVC}

L'aplicació segueix el patró Model-Vista-Controlador.
La responsabilitat de cada capa és clara:

\begin{table}[H]
\centering
\caption{Responsabilitats principals de cada capa.}
\label{tab:mvc}
\begin{tabular}{ll}
\toprule
Capa & Responsabilitat \\
\midrule
Model & Regles, tauler, moviments, resultat i agents \\
Vista & Swing, dibuix del tauler, controls i animacions \\
Controlador & Torns, modes, historial i connexió vista-model \\
\bottomrule
\end{tabular}
\end{table}

El model no depèn de Swing. Això permet que els agents utilitzin exactament les
mateixes regles que la interfície i facilita futurs canvis.
La vista tampoc modifica directament el joc: quan l'usuari
clica el tauler o un botó, genera un event cap al controlador. El controlador
decideix si el moviment és legal, aplica el model i demana a la vista que mostri
la transició.

\subsection{Decisions de Disseny}

La primera decisió rellevant és la immutabilitat del model. Encara que copiar un
tauler de $7\times7$ té un cost, aquest cost és petit i simplifica molt la cerca.
Cada branca del minimax pot aplicar moviments sense haver de desfer-los després.
Una alternativa seria tenir mètodes \texttt{doMove} i \texttt{undoMove}, però això
faria més fàcil introduir errors subtils amb rotacions i eliminacions.

La segona decisió és no introduir una condició d'empat per repetició d'estat. La
variant del joc pot generar cicles, però l'enunciat només defineix dues condicions
d'empat: victòria simultània i tauler ple sense guanyador. Per tant, la
implementació respecta estrictament aquestes condicions. Això implica que, en una
partida real, dos jugadors podrien repetir posicions; la interfície no ho talla
automàticament perquè no forma part de les regles demanades.

\subsection{Estructura de Paquets i Classes}

El projecte té un paquet base \texttt{com.serafinebot.p6} i subpaquets per a model,
controlador i vista. Aquesta organització evita dependències circulars i fa que
les responsabilitats es puguin localitzar ràpidament. La Taula~\ref{tab:classes}
resumeix les classes principals.

\begin{table}[H]
\centering
\caption{Classes principals de la pràctica.}
\label{tab:classes}
\begin{tabular}{ll}
\toprule
Classe & Funció \\
\midrule
\texttt{Main} & Punt d'entrada i arrencada Swing \\
\texttt{Board} & Tauler, gravetat i rotacions \\
\texttt{GameState} & Estat, torns i final de partida \\
\texttt{Move} & Representació d'una acció \\
\texttt{GameResult} & Resultat de la partida \\
\texttt{Agent} & Contracte comú dels agents \\
\texttt{RandomAgent} & Selecció aleatòria \\
\texttt{GreedyAgent} & Heurística local \\
\texttt{MinimaxAgent} & Cerca minimax amb poda \\
\texttt{GameController} & Orquestració MVC \\
\texttt{GameFrame} & Finestra principal \\
\texttt{BoardPanel} & Dibuix i animacions \\
\bottomrule
\end{tabular}
\end{table}

Les classes del model són deliberadament petites. \texttt{Player} només identifica
els dos jugadors i ofereix el mètode \texttt{opponent()}, que s'utilitza
constantment en el canvi de torn i en les heurístiques. \texttt{Cell} representa
el contingut d'una casella i inclou la conversió entre jugador i fitxa. Aquesta
separació evita haver d'emprar valors nuls dins del tauler: una cel·la buida és
\texttt{Cell.EMPTY}, no una absència d'objecte. \texttt{MoveType} enumera els
quatre tipus de moviment possibles i \texttt{Move} guarda les coordenades només
quan són necessàries.

Els enums \texttt{MatchType} i \texttt{AgentType} pertanyen a la vista perquè són
opcions de configuració de la interfície, no regles del joc. \texttt{MatchType}
distingeix humà contra humà, humà contra agent i agent contra agent; en canvi,
\texttt{AgentType} només indica quin agent concret s'ha de crear. Aquesta decisió
substitueix una enumeració combinatòria de modes per una configuració més
flexible: si s'afegeix un nou agent, no cal crear totes les combinacions possibles
de modes, només afegir un valor nou a \texttt{AgentType} i una branca de creació al
controlador.

\subsection{Flux Complet d'una Jugada}

El flux d'una jugada és el mateix per a humans i agents, amb l'única diferència de
qui proposa el moviment. En una jugada humana, \texttt{BoardPanel} rep el clic del
ratolí i calcula la fila i columna dins del tauler. Si la cel·la conté una fitxa
del jugador actual, envia \texttt{onRemove}; en cas contrari envia \texttt{onDrop}.
Els botons laterals de rotació envien \texttt{onRotateLeft} o
\texttt{onRotateRight}. En cap cas la vista modifica el model.

El controlador rep l'event i primer comprova que la partida no estigui ocupada amb
una animació o una decisió d'agent. Després comprova si el torn és realment humà.
Aquesta comprovació depèn del tipus de partida: en humà contra humà tots dos
jugadors són humans; en humà contra agent només ho és el vermell; en agent contra
agent no hi ha torns humans. Si el moviment és acceptable, el controlador crida
\texttt{state.apply(move)}.

El mètode \texttt{apply} no modifica l'estat existent. Construeix un tauler nou,
avalua el resultat i retorna un nou \texttt{GameState}. El controlador desa aquest
estat dins de l'historial, actualitza el cursor i demana a la vista que animi el
moviment. Mentre dura l'animació, la bandera \texttt{busy} impedeix nous clics. En
acabar, es tornen a habilitar els controls corresponents al nou torn. Aquest
detall és important perquè, si es permetessin entrades durant una animació, la
vista podria mostrar una transició que ja no correspon a l'estat real.

En una jugada d'agent, el procés és semblant però el moviment es calcula en segon
pla. \texttt{GameController} crea un \texttt{SwingWorker}, li passa una còpia de
l'estat actual i espera el moviment retornat per l'agent. Quan el treball acaba,
el moviment entra pel mateix camí que una jugada humana: s'aplica al model,
s'afegeix a l'historial i s'anima. Això evita duplicar regles entre jugadors
humans i automàtics.

\subsection{Model}

El model es divideix en tres blocs: representació física del tauler, estat complet
de partida i tipus auxiliars. Les subseccions següents descriuen aquests blocs amb
més detall; aquí només cal remarcar que totes les operacions de canvi d'estat
retornen objectes nous. Aquesta immutabilitat és la propietat que permet reutilitzar
el mateix model tant per a la partida real com per a les simulacions dels agents.

\subsubsection{Classe \texttt{Board}}

\texttt{Board} és la classe que conté més lògica física. El tauler és una matriu
\texttt{Cell[][]} de mida fixa, perquè l'enunciat fixa sempre $7\times7$. El
mètode \texttt{drop} comprova primer que la columna no estigui plena i després
recorre la columna des de baix cap amunt fins trobar la primera cel·la buida. Això
imita directament la caiguda d'una fitxa sense necessitat de simular frames
intermedis al model.

El mètode \texttt{remove} només accepta posicions ocupades per una fitxa del
jugador actual. Després de buidar la cel·la, crida \texttt{applyGravity}. Aquesta
funció compacta cada columna amb un índex \texttt{writeRow}: recorre la columna de
baix cap amunt i mou cada fitxa trobada a la posició lliure més baixa. En acabar,
totes les posicions que queden per damunt s'omplen amb \texttt{EMPTY}. Aquesta
tècnica és simple, no depèn de casos especials i garanteix que no quedin forats
entre fitxes.

Les rotacions es fan en dues fases. Primer es construeix una matriu nova amb la
transformació geomètrica corresponent. Per girar a l'esquerra, la posició
$(r,c)$ passa a $(N-1-c,r)$; per girar a la dreta passa a $(c,N-1-r)$. Després es
torna a aplicar gravetat, perquè una rotació pot deixar fitxes suspeses segons la
nova orientació. Aquesta separació entre rotació geomètrica i gravetat fa que la
implementació sigui fàcil de verificar.

La detecció de victòria també es concentra a \texttt{Board}. El mètode
\texttt{hasConnectFour} recorre totes les cel·les i, per cada fitxa del jugador
buscat, prova quatre direccions: dreta, avall, diagonal descendent i diagonal
ascendent. No cal provar les direccions contràries perquè trobarien les mateixes
línies duplicades. El cost és proporcional al nombre de cel·les multiplicat per
quatre direccions constants.

La correcció d'aquesta classe es pot expressar com una invariant de gravetat:
després de qualsevol operació pública, no hi ha cap fitxa amb una cel·la buida
immediatament davall dins de la mateixa columna. \texttt{drop} la preserva perquè
col·loca la fitxa directament a la posició lliure més baixa; \texttt{remove} i les
rotacions la restauren cridant \texttt{applyGravity}. Amb $N=7$, tant la gravetat
com la detecció de quatre en línia tenen cost $O(N^2)$ i són constants a efectes
pràctics.

\subsubsection{Classe \texttt{GameState}}

\texttt{GameState} és un record immutable amb tres camps: \texttt{board},
\texttt{currentPlayer} i \texttt{result}. La seva responsabilitat principal és
connectar les regles físiques del tauler amb la lògica de torns. Per això
\texttt{Board} sap eliminar una fitxa d'un jugador concret, però és
\texttt{GameState} qui decideix que el jugador concret és el jugador a qui toca
moure.

El mètode \texttt{legalMoves} genera els moviments en un ordre estable: primer les
insercions, després les eliminacions pròpies i finalment les dues rotacions. Aquest
ordre no és una regla del joc, però aporta determinisme als agents quan hi ha
empats de puntuació. També és important que les eliminacions només s'afegeixen si
el tauler no és ple, ja que l'enunciat especifica que amb totes les fitxes
col·locades la partida acaba en empat i no es pot continuar eliminant.

El mètode \texttt{evaluate} comprova sempre els dos jugadors després de cada
moviment. Aquesta decisió és necessària perquè una eliminació o una rotació pot
crear una línia per al jugador contrari. Si només es comprovàs el jugador que ha
mogut, el programa acceptaria partides incorrectes. La comprovació doble també
permet detectar correctament l'empat per victòria simultània.

La correcció de \texttt{GameState} depèn de l'ordre d'aquesta avaluació: primer es
calcula el tauler successor, després es comproven les línies guanyadores dels dos
jugadors, després l'empat per tauler ple i només si cap cas terminal apareix la
partida continua. La generació de moviments té cost $O(N+p)$, on $p$ és el nombre
de fitxes pròpies del jugador actual, perquè cal mirar les columnes disponibles i
les fitxes que es podrien eliminar.

\subsubsection{Resultats i Tipus Auxiliars}

\texttt{GameResult} és un record petit però útil: evita haver d'interpretar
manualment un \texttt{GameStatus} per saber si hi ha guanyador. En una victòria,
el camp \texttt{winner} conté el jugador; en empat o partida en curs és nul. Els
mètodes estàtics \texttt{win}, \texttt{draw} i \texttt{inProgress} fan que la
creació de resultats sigui clara i redueixen errors.

\texttt{Move} també és intencionadament simple. S'ha preferit un únic record amb
fàbriques estàtiques a una jerarquia de subclasses perquè els moviments són valors
petits que es creen massivament durant la cerca minimax. Aquesta representació
redueix la quantitat de classes, facilita imprimir moviments a l'estat de la
interfície i evita polimorfisme innecessari dins dels bucles de cerca.

\subsection{Vista i Interacció}

La vista està formada per \texttt{GameFrame}, que organitza la finestra, i
\texttt{BoardPanel}, que dibuixa el tauler. La decisió important és que cap de les
dues classes aplica regles de joc: només mostren l'estat rebut i transformen les
interaccions de l'usuari en events cap al controlador. Això evita que la lògica de
partida quedi duplicada entre Swing i el model.

\subsubsection{Classe \texttt{GameFrame}}

\texttt{GameFrame} és la finestra principal i implementa la interfície
\texttt{GameView}. La part superior conté el títol de l'aplicació i una breu
descripció de les accions disponibles. La zona central conté el tauler i dos
botons laterals de rotació. Aquesta disposició es va triar perquè la rotació és
una acció global sobre el tauler i visualment és més natural situar-la als costats
que dins del panell de configuració.

El panell lateral conté tres blocs. El primer és la configuració de partida: tipus
de partida i agents seleccionats. El segon és la reproducció: reproduir, aturar,
avançar i retrocedir. El tercer mostra els jugadors, el torn actual, el nombre de
fitxes col·locades i el darrer missatge del controlador. Els selectors d'agent es
deshabiliten dinàmicament segons el tipus de partida. En humà contra humà no fan
falta; en humà contra agent només s'empra l'agent groc; en agent contra agent
s'empren tots dos.

També s'ha implementat un botó personalitzat, \texttt{SolidButton}, perquè alguns

\textit{Look and Feel} de Swing, especialment en macOS, ignoren parcialment el
color de fons dels \texttt{JButton}. Aquesta classe pinta manualment un rectangle
arrodonit i canvia completament el color quan el botó està deshabilitat. Sense
això, alguns botons podien aparèixer amb text blanc sobre fons transparent.

\subsubsection{Classe \texttt{BoardPanel}}

\texttt{BoardPanel} és un component Swing dibuixat manualment. No utilitza una
graella de botons perquè el tauler necessita animacions i un aspecte visual més
proper al joc. El mètode \texttt{paintComponent} calcula la mida de cada cel·la
segons la mida disponible de la finestra, dibuixa el fons blau del tauler i pinta
cada fitxa com un cercle. Les fitxes eliminables es marquen amb una creu fosca,
de manera que el jugador veu quines peces pròpies pot retirar.

Per detectar clics s'utilitza \texttt{mousePressed} en lloc de
\texttt{mouseClicked}. Aquesta elecció fa que la resposta sigui més immediata:
\texttt{mouseClicked} només es dispara quan Swing reconeix una seqüència completa
de pressió i alliberament sense moviment, cosa que pot fallar si l'usuari mou
lleugerament el ratolí. Amb \texttt{mousePressed}, el panell respon tan aviat com
l'usuari prem dins del tauler.

L'animació de caiguda no duplica la lògica de gravetat. El panell rep l'estat
anterior i el posterior, compara la columna afectada i troba la fila on ha aparegut
la nova fitxa. A partir d'aquí només interpola la coordenada vertical amb una
funció cúbica \textit{ease-out}. Això dóna una caiguda ràpida al principi i més
suau al final. En les rotacions, el panell guarda temporalment l'estat anterior i
el dibuixa amb una \texttt{AffineTransform}; quan acaba el temporitzador, torna a
mostrar l'estat final normal.

\subsection{Controlador i Historial}

\texttt{GameController} és el punt central de coordinació. Manté l'estat actual,
els agents seleccionats, el tipus de partida i un historial de posicions. També
manté un cursor dins d'aquest historial per implementar els botons \textit{Següent}
i \textit{Anterior}. Si l'usuari retrocedeix i després fa una jugada nova, el
controlador elimina el futur antic de l'historial, de la mateixa manera que un
editor elimina el camí de refer després d'una modificació.

Les decisions dels agents s'executen dins d'un \texttt{SwingWorker}. Aquesta
decisió és necessària perquè el minimax pot explorar milers de nodes i, si
s'executàs directament al fil d'events de Swing, congelaria la interfície. En mode
agent contra agent, un \texttt{Timer} de Swing crida periòdicament el següent
moviment amb un retard visible, cosa que permet observar la partida en temps real.

El controlador també centralitza l'activació i desactivació de controls. Aquesta
responsabilitat no es deixa a la vista perquè depèn de l'estat lògic del joc. Per
exemple, els botons de rotació han d'estar actius en un torn humà, però no mentre
un agent està pensant ni durant una animació. De la mateixa manera, el botó
\textit{Reproduir} només té sentit en mode agent contra agent i quan la partida no
ha acabat. Aquesta política queda concentrada al mètode \texttt{refresh}, que
actualitza l'estat visual després de cada canvi.

La llista \texttt{history} desa tots els \texttt{GameState} visitats i la llista
\texttt{moves} desa el moviment que separa cada parella d'estats consecutius. El
cursor indica quina posició de l'historial s'està mostrant. Quan l'usuari prem
\textit{Anterior}, el controlador redueix el cursor i mostra directament l'estat
anterior. Quan prem \textit{Següent}, si el moviment ja existeix a l'historial, el
controlador avança sense animació perquè és navegació històrica; si no existeix,
demana un moviment nou a l'agent corresponent. Aquest comportament separa clarament
la reproducció d'una partida ja calculada de la generació d'una jugada nova.

%% ================================================================
\section{Agents Implementats}

\subsection{Agent Aleatori}

\texttt{RandomAgent} és l'agent base. Obté la llista de moviments legals i en tria
un uniformement a l'atzar. No diferencia entre inserir, eliminar o rotar, ni mira
si el moviment guanya o perd. El seu valor és de control: serveix per comprovar
que el motor de joc accepta moviments legals encara que no hi hagi cap criteri
estratègic darrere la decisió.

Aquest agent també és útil durant el desenvolupament. Com que no té cap lògica
estratègica, si una partida aleatòria produeix un moviment il·legal o una excepció,
el problema gairebé sempre és al model o al controlador, no a l'agent. A més,
permet provar ràpidament les animacions i els modes agent contra agent sense
esperar el cost del minimax.

La complexitat temporal d'una decisió és $O(|M(S)|)$ per generar els moviments i
$O(1)$ per seleccionar-ne un. La qualitat de joc és baixa però el temps de decisió
és pràcticament constant per a l'usuari.

\subsection{Agent Greedy}

\texttt{GreedyAgent} és un agent d'una sola capa. Per cada moviment legal, aplica
el moviment sobre una còpia immutable de l'estat i valora només la posició
resultant. Si el moviment produeix una victòria pròpia, rep una puntuació molt
alta; si produeix una victòria de l'oponent, rep una puntuació molt baixa; si no
és terminal, es calcula una suma simple de material i control del centre.

Aquest agent és més fort que l'aleatori perquè reconeix victòries immediates i
prefereix posicions amb més presència pròpia. Tot i això, no modela la resposta de
l'oponent. Pot triar una jugada aparentment bona que permeti al contrari guanyar
immediatament al torn següent. Per aquest motiu és una bona referència intermèdia
entre l'agent aleatori i el minimax.

Si $b$ és el nombre de moviments legals i $C$ el cost de valorar un tauler, el
cost és $O(b\cdot C)$. En aquesta implementació $C$ és proporcional a les 49
cel·les del tauler, de manera que el cost pràctic és molt baix.

La funció de puntuació greedy no utilitza finestres de quatre cel·les com el
minimax. S'ha mantingut més senzilla expressament perquè el seu paper és ser un
intermedi entre l'aleatori i la cerca adversarial. Si el greedy tengués una
heurística gairebé igual al minimax, seria més difícil distingir si la millora de
rendiment ve de la cerca en profunditat o només de la funció d'avaluació. Amb
material i centre, el greedy té criteri però continua essent clarament local.

\subsection{Agent Minimax}

\texttt{MinimaxAgent} explora l'arbre de joc fins a una profunditat fixa. En cada
nivell alterna maximització i minimització segons el jugador que ha de moure. Les
posicions terminals obtenen valors extrems i les posicions no terminals al límit
de profunditat es valoren amb la funció heurística descrita anteriorment.

La profunditat per defecte és de 3 nivells de moviment. En la terminologia de
jocs, un nivell d'aquest tipus s'anomena sovint \textit{ply}: és una jugada
individual d'un sol jugador, no una ronda completa de dos torns. Per tant, una
profunditat 3 vol dir que l'agent considera la seva jugada, la resposta de
l'oponent i una jugada pròpia addicional. Aquesta profunditat és un compromís
pràctic: prou gran per detectar moltes amenaces immediates, però prou petita per
respondre en una interfície interactiva. Augmentar-la milloraria la qualitat en
alguns casos, però el factor de ramificació de la variant fa que el cost creixi
ràpidament.

El pseudocodi següent mostra tant la selecció del moviment inicial com la funció
recursiva. Els paràmetres $\alpha$ i $\beta$ són els límits de poda: quan es
creuen, la branca actual no pot canviar la decisió final i s'abandona.

\begin{algorithm}[H]
\caption{Minimax amb poda alfa-beta}
\begin{algorithmic}[1]
\STATE \textbf{Selecció inicial del moviment}
\STATE $millorValor \leftarrow -\infty$, $millorMoviment \leftarrow$ primer moviment
\STATE $\alpha \leftarrow -\infty$, $\beta \leftarrow +\infty$
\FORALL{$m \in$ moviments ordenats de $s$}
    \STATE $s' \leftarrow \operatorname{apply}(s,m)$
    \STATE $v \leftarrow V(s', d-1, \alpha, \beta, jugadorMax)$
    \IF{$v > millorValor$}
        \STATE $millorValor \leftarrow v$, $millorMoviment \leftarrow m$
    \ENDIF
    \STATE $\alpha \leftarrow \max(\alpha, millorValor)$
\ENDFOR
\STATE \textbf{return} $millorMoviment$
\STATE
\STATE \textbf{Funció recursiva} $V(s,d,\alpha,\beta,jugadorMax)$
\IF{$d=0$ o $s$ és terminal}
    \STATE \textbf{return} $h(s)$
\ENDIF
\STATE $maximitza \leftarrow$ el jugador de $s$ és $jugadorMax$
\IF{$maximitza$}
    \STATE $valor \leftarrow -\infty$
\ELSE
    \STATE $valor \leftarrow +\infty$
\ENDIF
\FORALL{$m \in$ moviments ordenats de $s$}
    \STATE $s_m \leftarrow \operatorname{apply}(s,m)$
    \STATE $fill \leftarrow V(s_m,d-1,\alpha,\beta,jugadorMax)$
    \IF{$maximitza$}
        \STATE $valor \leftarrow \max(valor, fill)$
        \STATE $\alpha \leftarrow \max(\alpha, valor)$
    \ELSE
        \STATE $valor \leftarrow \min(valor, fill)$
        \STATE $\beta \leftarrow \min(\beta, valor)$
    \ENDIF
    \IF{$\alpha \geq \beta$}
        \STATE interrompre el bucle
    \ENDIF
\ENDFOR
\STATE \textbf{return} $valor$
\end{algorithmic}
\end{algorithm}

Sense poda, el cost temporal és $O(b^d)$, on $d$ és la profunditat~\cite{ref_gfg_minimax}. Amb alfa-beta,
el pitjor cas continua essent $O(b^d)$, però amb una bona ordenació de moviments el
nombre de nodes visitats es redueix molt. En el cas ideal clàssic, la complexitat
efectiva s'aproxima a $O(b^{d/2})$~\cite{ref_gfg_alphabeta}. Per això el programa desa
estadístiques de nodes explorats, podes i temps. En aquesta memòria s'utilitzen
per justificar el cost d'una decisió concreta, no com a base d'una comparativa
estadística.

\subsection{Ordenació de Moviments i Estadístiques}

Abans d'explorar successors, \texttt{MinimaxAgent} els ordena amb una puntuació
ràpida. Aquesta puntuació no és la mateixa que la heurística completa: només serveix
per decidir l'ordre de visita. Les victòries immediates es posen al principi,
després es prefereixen insercions properes al centre i les altres accions queden
amb menor prioritat. L'objectiu és que alfa-beta trobi aviat bons límits i pugui
tallar més branques.

L'agent guarda tres mesures de la darrera cerca: nodes visitats, podes realitzades
i temps de decisió en mil·lisegons. Aquestes dades es mostren al missatge d'estat
quan juga el minimax. No s'han emprat per construir una taula experimental, però
sí per comprovar que la profunditat triada manté un temps de resposta adequat per
a una interfície interactiva.

\subsection{Discussió Qualitativa dels Agents}

Els tres agents representen tres nivells d'intel·ligència artificial. L'aleatori
no utilitza cap informació del joc i, per tant, és una línia base mínima. El
greedy incorpora una noció de posició però només mira un moviment. El minimax
incorpora oposició explícita: assumeix que el contrari també intentarà obtenir el
millor resultat possible. Aquesta diferència és essencial en jocs adversarials,
perquè una jugada no s'ha de valorar només pel que crea ara, sinó també pel que
permet respondre.

En aquesta variant concreta, la diferència entre greedy i minimax és especialment
visible amb rotacions i eliminacions. Una rotació pot semblar dolenta localment
perquè no augmenta el material ni ocupa el centre, però pot preparar o impedir una
línia diagonal després de la resposta del contrari. El greedy tendeix a infravalorar
aquests moviments; el minimax, en canvi, pot acceptar una jugada amb puntuació
local modesta si les respostes futures són favorables.

%% ================================================================
\section{Resultats Funcionals i Discussió}
\label{sec:resultats}

En aquesta pràctica no s'ha fet una campanya experimental amb moltes partides ni
s'han recollit dades estadístiques de victòries, temps mitjans o percentatges de
poda. Per tant, aquest apartat no presenta una comparativa quantitativa entre
agents. Els resultats que es poden defensar són funcionals: el programa implementa
les regles exigides, permet jugar amb diferents tipus de participants i mostra el
cost de cada decisió minimax quan aquest agent intervé.

La Taula~\ref{tab:functional-results} resumeix les comprovacions principals de la
implementació. No són dades experimentals comparatives, sinó casos de validació
que confirmen que les parts essencials del joc responen segons les regles de
l'enunciat.

\begin{table}[H]
\centering
\caption{Resultats funcionals validats.}
\label{tab:functional-results}
\begin{tabular}{p{0.35\columnwidth}p{0.53\columnwidth}}
\toprule
Funcionalitat & Resultat comprovat \\
\midrule
Inserció & La fitxa cau fins a la cel·la lliure més baixa. \\
Eliminació & Només es poden llevar fitxes pròpies i s'aplica gravetat. \\
Rotació & El tauler gira 90 graus i es compacta segons la nova orientació. \\
Final de partida & Es detecten victòria, empat simultani i tauler ple. \\
Agents & Els tres agents generen sempre moviments legals. \\
Historial & Es pot revisar una partida automàtica endavant i enrere. \\
\bottomrule
\end{tabular}
\end{table}

La discussió dels agents s'ha de llegir en aquest mateix sentit qualitatiu. L'agent
aleatori serveix per comprovar que el motor accepta qualsevol moviment legal sense
coneixement estratègic. L'agent greedy comprova que una heurística local pot
aprofitar oportunitats immediates, però també mostra la limitació de no anticipar
la resposta rival. El minimax és l'agent que satisfà el requisit algorísmic central:
explora diversos nivells, assigna valors als estats futurs i talla branques amb
poda alfa-beta.

Les estadístiques mostrades pel controlador per a \texttt{MinimaxAgent} tenen una
funció diagnòstica. El nombre de nodes, el nombre de podes i el temps de decisió
permeten veure si una jugada concreta ha requerit més cerca que una altra, però no
s'han agregat en una taula de rendiment. Per obtenir una comparativa empírica real,
caldria executar moltes partides automàtiques per parella d'agents, alternar el
color inicial i calcular mitjanes i percentatges. Això queda fora de l'abast
implementat aquí.

%% ================================================================
\section{Conclusions}

S'ha desenvolupat una aplicació Java completa per jugar a la variant proposada de
Connecta 4. El model implementa les regles essencials: inserció amb gravetat,
eliminació de fitxes pròpies, rotació del tauler, detecció de quatre en línia,
empat simultani i empat per tauler ple. La separació MVC facilita mantenir les
regles independents de la interfície i permet reutilitzar el mateix model per
humans, agents i futures comparatives.

Pel que fa als agents, s'han implementat tres nivells de complexitat. L'agent
aleatori serveix de referència mínima, el greedy introdueix una heurística local i
el minimax amb poda alfa-beta satisfà el requisit de cerca adversarial amb Branch
and Bound. La funció heurística del minimax combina finestres de quatre cel·les,
control del centre i material, amb pesos que prioritzen bloquejar amenaces
immediates.

Durant les proves manuals s'ha observat que és molt difícil guanyar contra el
minimax, però no impossible. Aquest comportament és coherent amb el fet que
l'agent no resol el joc complet: utilitza una profunditat limitada i una heurística
aproximada, de manera que pot deixar escapar posicions favorables quan la seqüència
crítica queda més enllà de l'horitzó de cerca o no és valorada prou bé.

Com a extensions naturals, es podrien afegir profunditats configurables per al
minimax i una comparació explícita entre minimax amb i sense ordenació de
moviments.

\begin{thebibliography}{9}

\bibitem{ref_gfg_minimax}
GeeksforGeeks, ``Minimax Algorithm in Game Theory | Set 1 (Introduction),''
2026. [Online]. Available:
\url{https://www.geeksforgeeks.org/dsa/minimax-algorithm-in-game-theory-set-1-introduction/}

\bibitem{ref_gfg_alphabeta}
GeeksforGeeks, ``Minimax Algorithm in Game Theory | Set 4 (Alpha-Beta
Pruning),'' 2023. [Online]. Available:
\url{https://www.geeksforgeeks.org/dsa/minimax-algorithm-in-game-theory-set-4-alpha-beta-pruning/}

\end{thebibliography}

\end{document}
